%=====================================================================
%  backend.tex : MMSS Graph Theory report engine
%
%  DO NOT edit this file for day-to-day reports. Edit main.tex.
%  Touch this file only when you want a NEW macro or a style change.
%  (Makers note : This was really hard to make and I strongly suggest 
%  you dont touch it T-T....)
%  Loaded by main.tex via %=====================================================================
%  backend.tex : MMSS Graph Theory report engine
%
%  DO NOT edit this file for day-to-day reports. Edit main.tex.
%  Touch this file only when you want a NEW macro or a style change.
%  (Makers note : This was really hard to make and I strongly suggest 
%  you dont touch it T-T....)
%  Loaded by main.tex via %=====================================================================
%  backend.tex : MMSS Graph Theory report engine
%
%  DO NOT edit this file for day-to-day reports. Edit main.tex.
%  Touch this file only when you want a NEW macro or a style change.
%  (Makers note : This was really hard to make and I strongly suggest 
%  you dont touch it T-T....)
%  Loaded by main.tex via %=====================================================================
%  backend.tex : MMSS Graph Theory report engine
%
%  DO NOT edit this file for day-to-day reports. Edit main.tex.
%  Touch this file only when you want a NEW macro or a style change.
%  (Makers note : This was really hard to make and I strongly suggest 
%  you dont touch it T-T....)
%  Loaded by main.tex via \input{backend} (after \documentclass).
%
%  Expects (optionally) these to be defined in main.tex BEFORE \input:
%    \teamname    \reportdate    \logopath
%  Fallbacks are provided below, so nothing breaks if you forget.
%
%  ---------------------------------------------------------------
%  GRAPH TOOLKIT CHEAT SHEET  (all calls go INSIDE graphfig/tikzpicture)
%  ---------------------------------------------------------------
%    \gcycle{n}{radius}{prefix}      n vertices on a circle, joined as a cycle
%    \gcyclen{n}{radius}{prefix}     same, vertices labelled 1..n
%    \gcomplete{n}{radius}{prefix}   complete graph K_n on a circle
%    \gbipartite{m}{n}{gap}          complete bipartite K_{m,n}
%                                      handles: left bl1..blm, right br1..brn
%    \gpath{n}{spacing}{prefix}      path on n vertices, left to right
%    \gpetersen                      Petersen graph  (\gpeterson kept as alias)
%                                      handles: outer po0..po4, inner pin0..pin4
%    \gcube{n}                       n x n x n grid cube, oblique 2D projection
%    \gqgon{n}, \gpancake{n}         suprise hidden functions, only use after you 
%                                      learned what it is. 
%
%  Named handles stay live for the rest of the picture, so after a call:
%    \draw[edge,densely dashed] (po0)--(pin2);   or   \node[hv] at (bl1){};
%
%  Vertex styles:  v (filled dot) . ov (open ring) . hv (highlighted hub)
%  Edge style:     edge
%
%  Figure wrapper: \begin{graphfig}[scale]{Label} ... \end{graphfig}
%    {Text} -> centred, titled  [Graph Text]
%    {NUM}  -> centred, auto-numbered  [Graph 1], [Graph 2], ...
%    {}     -> bare, no title, no centring (use to line pictures up side by side)
%=====================================================================

\usepackage[margin=1in,headheight=30pt,headsep=16pt]{geometry}
\usepackage{amsmath,amssymb,amsfonts,amsthm}
\usepackage{enumitem}
\usepackage{fancyhdr}
\usepackage{graphicx}
\usepackage{tikz}
\usepackage{xcolor}
\usepackage{hyperref}
\usetikzlibrary{arrows.meta,calc}


\providecommand{\teamname}{[GROUP NAME]}
\providecommand{\reportdate}{\today}
\providecommand{\logopath}{assets/logo.png}

\pagestyle{fancy}
\fancyhf{}
\renewcommand{\headrulewidth}{0.4pt}
\renewcommand{\footrulewidth}{0pt}
\fancyhead[L]{\raisebox{2pt}{\small\textsc{\teamname}}}
\fancyhead[R]{\expandafter\IfFileExists\expandafter{\logopath}
  {\includegraphics[height=10pt]{\logopath}}{}}
\fancyfoot[C]{\small\thepage}

\fancypagestyle{firstpage}{
  \fancyhf{}
  \renewcommand{\headrulewidth}{0pt}
  \fancyfoot[C]{\small\thepage}
}

\setlength{\parindent}{0pt}
\setlength{\parskip}{6pt}

\newcommand{\divider}{\begin{center}\rule{0.6\textwidth}{0.4pt}\end{center}}

\newtheoremstyle{solstyle}{6pt}{6pt}{}{}{\bfseries}{.}{ }{}
\theoremstyle{solstyle}
\newtheorem*{solution}{Solution}

\newcommand{\WLoG}{WLoG}
\newcommand{\Cn}[1]{C_{#1}}
\newcommand{\Kn}[1]{K_{#1}}
\newcommand{\Kt}[2]{K_{#1,#2}}

\tikzset{
  v/.style ={circle, fill=black, inner sep=1.6pt},
  ov/.style={circle, draw=black, fill=white, line width=0.6pt, inner sep=1.6pt},
  hv/.style={circle, fill=black!55, draw=black, line width=0.6pt, inner sep=2.6pt},
  edge/.style={line width=0.6pt},
}


\newcommand{\gcycle}[3]{
  \foreach \gi in {1,...,#1}{\node[v] (#3\gi) at ({90-(\gi-1)*360/#1}:#2){};}
  \foreach \gi in {1,...,#1}{\pgfmathtruncatemacro{\gj}{mod(\gi,#1)+1}\draw[edge] (#3\gi)--(#3\gj);}
}

\newcommand{\gcyclen}[3]{
  \foreach \gi in {1,...,#1}{
    \pgfmathsetmacro{\gang}{90-(\gi-1)*360/#1}
    \node[v,label={\gang:\gi}] (#3\gi) at (\gang:#2){};}
  \foreach \gi in {1,...,#1}{\pgfmathtruncatemacro{\gj}{mod(\gi,#1)+1}\draw[edge] (#3\gi)--(#3\gj);}
}

\newcommand{\gcomplete}[3]{
  \foreach \gi in {1,...,#1}{\node[v] (#3\gi) at ({90-(\gi-1)*360/#1}:#2){};}
  \foreach \gi in {1,...,#1}{\foreach \gj in {1,...,#1}{\ifnum\gi<\gj\relax\draw[edge] (#3\gi)--(#3\gj);\fi}}
}

\newcommand{\gbipartite}[3]{
  \foreach \gi in {1,...,#1}{\node[v] (bl\gi) at (0,{(#1+1)/2-\gi}){};}
  \foreach \gj in {1,...,#2}{\node[v] (br\gj) at (#3,{(#2+1)/2-\gj}){};}
  \foreach \gi in {1,...,#1}{\foreach \gj in {1,...,#2}{\draw[edge] (bl\gi)--(br\gj);}}
}

\newcommand{\gpath}[3]{
  \foreach \gi in {1,...,#1}{\node[v] (#3\gi) at ({(\gi-1)*#2},0){};}
  \pgfmathtruncatemacro{\gplast}{#1-1}
  \foreach \gi in {1,...,\gplast}{\pgfmathtruncatemacro{\gpk}{\gi+1}\draw[edge] (#3\gi)--(#3\gpk);}
}


\newcommand{\gpetersen}{
  \foreach \gp in {0,...,4}{
    \node[v] (po\gp) at ({90+\gp*72}:2){};
    \node[v] (pin\gp) at ({90+\gp*72}:1){};}
  \foreach \gp in {0,...,4}{
    \pgfmathtruncatemacro{\gpj}{mod(\gp+1,5)}
    \pgfmathtruncatemacro{\gpk}{mod(\gp+2,5)}
    \draw[edge] (po\gp)--(po\gpj);
    \draw[edge] (po\gp)--(pin\gp);
    \draw[edge] (pin\gp)--(pin\gpk);
  }
}
\newcommand{\gpeterson}{\gpetersen}

\newcommand{\gcube}[1]{
  \begin{scope}[x={(1cm,0cm)}, y={(0cm,1cm)}, z={(0.5cm,0.45cm)}]
    \foreach \ci in {0,...,#1}{\draw[edge] (\ci,0,0)--(\ci,#1,0);}
    \foreach \cj in {0,...,#1}{\draw[edge] (0,\cj,0)--(#1,\cj,0);}
    \foreach \ci in {0,...,#1}{\draw[edge] (\ci,#1,0)--(\ci,#1,#1);}
    \foreach \ck in {0,...,#1}{\draw[edge] (0,#1,\ck)--(#1,#1,\ck);}
    \foreach \cj in {0,...,#1}{\draw[edge] (#1,\cj,0)--(#1,\cj,#1);}
    \foreach \ck in {0,...,#1}{\draw[edge] (#1,0,\ck)--(#1,#1,\ck);}
    \foreach \ci in {0,...,#1}{\foreach \cj in {0,...,#1}{\node[v] at (\ci,\cj,0){};}}
    \foreach \ci in {0,...,#1}{\foreach \ck in {0,...,#1}{\node[v] at (\ci,#1,\ck){};}}
    \foreach \cj in {0,...,#1}{\foreach \ck in {0,...,#1}{\node[v] at (#1,\cj,\ck){};}}
  \end{scope}
}

\newcount\pgPlaced
\newcount\pgNf
\newcount\pgTt
\newcommand{\gpancake}[2]{
  \pgNf=1 \pgTt=1
  \loop\ifnum\pgTt<#1 \advance\pgTt by1\relax \multiply\pgNf by\pgTt \repeat
  \edef\pgNfact{\the\pgNf}
  \pgfmathtruncatemacro{\pgTotm}{#1^#1-1}
  \pgPlaced=0
  \foreach \pgI in {0,...,\pgTotm}{
    \gdef\pgPerm{}\xdef\pgUsed{0}\xdef\pgValid{1}
    \foreach \pos in {1,...,#1}{
      \pgfmathtruncatemacro{\pgDig}{Mod(floor(\pgI/(#1^(#1-\pos))),#1)}
      \pgfmathtruncatemacro{\pgBit}{Mod(floor(\pgUsed/(2^\pgDig)),2)}
      \ifnum\pgBit=1 \xdef\pgValid{0}
      \else
        \pgfmathtruncatemacro{\pgNewUsed}{\pgUsed+2^\pgDig}\xdef\pgUsed{\pgNewUsed}
        \pgfmathtruncatemacro{\pgD}{\pgDig+1}\xdef\pgPerm{\pgPerm\pgD}
      \fi}
    \ifnum\pgValid=1
      \pgfmathsetmacro{\pgAng}{90-(\the\pgPlaced)*360/\pgNfact}
      \node[v,label={\pgAng:{\tiny\pgPerm}}] (p\pgPerm) at (\pgAng:#2){};
      \global\advance\pgPlaced by1\relax
    \fi}
  \foreach \pgI in {0,...,\pgTotm}{
    \gdef\pgPerm{}\xdef\pgUsed{0}\xdef\pgValid{1}
    \foreach \pos in {1,...,#1}{
      \pgfmathtruncatemacro{\pgDig}{Mod(floor(\pgI/(#1^(#1-\pos))),#1)}
      \pgfmathtruncatemacro{\pgBit}{Mod(floor(\pgUsed/(2^\pgDig)),2)}
      \ifnum\pgBit=1 \xdef\pgValid{0}
      \else
        \pgfmathtruncatemacro{\pgNewUsed}{\pgUsed+2^\pgDig}\xdef\pgUsed{\pgNewUsed}
        \pgfmathtruncatemacro{\pgD}{\pgDig+1}\xdef\pgPerm{\pgPerm\pgD}
        \expandafter\xdef\csname pgDig\pos\endcsname{\pgD}
      \fi}
    \ifnum\pgValid=1
      \foreach \k in {2,...,#1}{
        \xdef\pgNb{}
        \foreach \j in {1,...,\k}{\xdef\pgNb{\csname pgDig\j\endcsname\pgNb}}
        \ifnum\k<#1
          \foreach \j in {\the\numexpr\k+1\relax,...,#1}{\xdef\pgNb{\pgNb\csname pgDig\j\endcsname}}
        \fi
        \draw[edge] (p\pgPerm)--(p\pgNb);
      }
    \fi}
}




\newcommand{\gqgon}[2]{
  \pgfmathtruncatemacro{\Nm}{2^#1-1}
  \pgfmathtruncatemacro{\nm}{#1-1}
  \foreach \i in {0,...,\Nm}{
    \pgfmathsetmacro{\ang}{90-\i*360/(\Nm+1)}
    \xdef\lbl{}
    \foreach \k in {\nm,...,0}{
      \pgfmathtruncatemacro{\bit}{mod(floor(\i/(2^\k)),2)}
      \xdef\lbl{\lbl\bit}
    }
    \node[v,label={\ang:\scriptsize $\lbl$}] (q\i) at (\ang:#2){};
  }
  \foreach \i in {0,...,\Nm}{
    \foreach \k in {0,...,\nm}{
      \pgfmathtruncatemacro{\bit}{mod(floor(\i/(2^\k)),2)}
      \ifnum\bit=0
        \pgfmathtruncatemacro{\j}{\i+2^\k}
        \draw[edge] (q\i)--(q\j);
      \fi
    }
  }
}




\newcommand{\newinstructor}[2]{
  \definecolor{#2color}{HTML}{#1}
  \expandafter\newcommand\csname #2\endcsname[1]{
    {\color{#2color}\textbf{#2} : ##1}}}

\newcounter{graphnum}
\newif\ifgfbare
\def\gfEmpty{}
\def\gfNUM{NUM}
\newenvironment{graphfig}[2][1]{
  \def\gftmp{#2}
  \ifx\gftmp\gfEmpty
    \gfbaretrue
  \else
    \gfbarefalse
    \ifx\gftmp\gfNUM
      \stepcounter{graphnum}\edef\gflabel{\arabic{graphnum}}
    \else
      \def\gflabel{#2}
    \fi
  \fi
  \ifgfbare\else\begin{center}\fi
  \begin{tikzpicture}[scale=#1]
}{
  \end{tikzpicture}
  \ifgfbare\else\par\vspace{3pt}\textbf{[Graph \gflabel]}\end{center}\fi
}

\newcommand{\glabel}[1]{\begin{center}\textbf{[Graph #1]}\end{center}} (after \documentclass).
%
%  Expects (optionally) these to be defined in main.tex BEFORE \input:
%    \teamname    \reportdate    \logopath
%  Fallbacks are provided below, so nothing breaks if you forget.
%
%  ---------------------------------------------------------------
%  GRAPH TOOLKIT CHEAT SHEET  (all calls go INSIDE graphfig/tikzpicture)
%  ---------------------------------------------------------------
%    \gcycle{n}{radius}{prefix}      n vertices on a circle, joined as a cycle
%    \gcyclen{n}{radius}{prefix}     same, vertices labelled 1..n
%    \gcomplete{n}{radius}{prefix}   complete graph K_n on a circle
%    \gbipartite{m}{n}{gap}          complete bipartite K_{m,n}
%                                      handles: left bl1..blm, right br1..brn
%    \gpath{n}{spacing}{prefix}      path on n vertices, left to right
%    \gpetersen                      Petersen graph  (\gpeterson kept as alias)
%                                      handles: outer po0..po4, inner pin0..pin4
%    \gcube{n}                       n x n x n grid cube, oblique 2D projection
%    \gqgon{n}, \gpancake{n}         suprise hidden functions, only use after you 
%                                      learned what it is. 
%
%  Named handles stay live for the rest of the picture, so after a call:
%    \draw[edge,densely dashed] (po0)--(pin2);   or   \node[hv] at (bl1){};
%
%  Vertex styles:  v (filled dot) . ov (open ring) . hv (highlighted hub)
%  Edge style:     edge
%
%  Figure wrapper: \begin{graphfig}[scale]{Label} ... \end{graphfig}
%    {Text} -> centred, titled  [Graph Text]
%    {NUM}  -> centred, auto-numbered  [Graph 1], [Graph 2], ...
%    {}     -> bare, no title, no centring (use to line pictures up side by side)
%=====================================================================

\usepackage[margin=1in,headheight=30pt,headsep=16pt]{geometry}
\usepackage{amsmath,amssymb,amsfonts,amsthm}
\usepackage{enumitem}
\usepackage{fancyhdr}
\usepackage{graphicx}
\usepackage{tikz}
\usepackage{xcolor}
\usepackage{hyperref}
\usetikzlibrary{arrows.meta,calc}


\providecommand{\teamname}{[GROUP NAME]}
\providecommand{\reportdate}{\today}
\providecommand{\logopath}{assets/logo.png}

\pagestyle{fancy}
\fancyhf{}
\renewcommand{\headrulewidth}{0.4pt}
\renewcommand{\footrulewidth}{0pt}
\fancyhead[L]{\raisebox{2pt}{\small\textsc{\teamname}}}
\fancyhead[R]{\expandafter\IfFileExists\expandafter{\logopath}
  {\includegraphics[height=10pt]{\logopath}}{}}
\fancyfoot[C]{\small\thepage}

\fancypagestyle{firstpage}{
  \fancyhf{}
  \renewcommand{\headrulewidth}{0pt}
  \fancyfoot[C]{\small\thepage}
}

\setlength{\parindent}{0pt}
\setlength{\parskip}{6pt}

\newcommand{\divider}{\begin{center}\rule{0.6\textwidth}{0.4pt}\end{center}}

\newtheoremstyle{solstyle}{6pt}{6pt}{}{}{\bfseries}{.}{ }{}
\theoremstyle{solstyle}
\newtheorem*{solution}{Solution}

\newcommand{\WLoG}{WLoG}
\newcommand{\Cn}[1]{C_{#1}}
\newcommand{\Kn}[1]{K_{#1}}
\newcommand{\Kt}[2]{K_{#1,#2}}

\tikzset{
  v/.style ={circle, fill=black, inner sep=1.6pt},
  ov/.style={circle, draw=black, fill=white, line width=0.6pt, inner sep=1.6pt},
  hv/.style={circle, fill=black!55, draw=black, line width=0.6pt, inner sep=2.6pt},
  edge/.style={line width=0.6pt},
}


\newcommand{\gcycle}[3]{
  \foreach \gi in {1,...,#1}{\node[v] (#3\gi) at ({90-(\gi-1)*360/#1}:#2){};}
  \foreach \gi in {1,...,#1}{\pgfmathtruncatemacro{\gj}{mod(\gi,#1)+1}\draw[edge] (#3\gi)--(#3\gj);}
}

\newcommand{\gcyclen}[3]{
  \foreach \gi in {1,...,#1}{
    \pgfmathsetmacro{\gang}{90-(\gi-1)*360/#1}
    \node[v,label={\gang:\gi}] (#3\gi) at (\gang:#2){};}
  \foreach \gi in {1,...,#1}{\pgfmathtruncatemacro{\gj}{mod(\gi,#1)+1}\draw[edge] (#3\gi)--(#3\gj);}
}

\newcommand{\gcomplete}[3]{
  \foreach \gi in {1,...,#1}{\node[v] (#3\gi) at ({90-(\gi-1)*360/#1}:#2){};}
  \foreach \gi in {1,...,#1}{\foreach \gj in {1,...,#1}{\ifnum\gi<\gj\relax\draw[edge] (#3\gi)--(#3\gj);\fi}}
}

\newcommand{\gbipartite}[3]{
  \foreach \gi in {1,...,#1}{\node[v] (bl\gi) at (0,{(#1+1)/2-\gi}){};}
  \foreach \gj in {1,...,#2}{\node[v] (br\gj) at (#3,{(#2+1)/2-\gj}){};}
  \foreach \gi in {1,...,#1}{\foreach \gj in {1,...,#2}{\draw[edge] (bl\gi)--(br\gj);}}
}

\newcommand{\gpath}[3]{
  \foreach \gi in {1,...,#1}{\node[v] (#3\gi) at ({(\gi-1)*#2},0){};}
  \pgfmathtruncatemacro{\gplast}{#1-1}
  \foreach \gi in {1,...,\gplast}{\pgfmathtruncatemacro{\gpk}{\gi+1}\draw[edge] (#3\gi)--(#3\gpk);}
}


\newcommand{\gpetersen}{
  \foreach \gp in {0,...,4}{
    \node[v] (po\gp) at ({90+\gp*72}:2){};
    \node[v] (pin\gp) at ({90+\gp*72}:1){};}
  \foreach \gp in {0,...,4}{
    \pgfmathtruncatemacro{\gpj}{mod(\gp+1,5)}
    \pgfmathtruncatemacro{\gpk}{mod(\gp+2,5)}
    \draw[edge] (po\gp)--(po\gpj);
    \draw[edge] (po\gp)--(pin\gp);
    \draw[edge] (pin\gp)--(pin\gpk);
  }
}
\newcommand{\gpeterson}{\gpetersen}

\newcommand{\gcube}[1]{
  \begin{scope}[x={(1cm,0cm)}, y={(0cm,1cm)}, z={(0.5cm,0.45cm)}]
    \foreach \ci in {0,...,#1}{\draw[edge] (\ci,0,0)--(\ci,#1,0);}
    \foreach \cj in {0,...,#1}{\draw[edge] (0,\cj,0)--(#1,\cj,0);}
    \foreach \ci in {0,...,#1}{\draw[edge] (\ci,#1,0)--(\ci,#1,#1);}
    \foreach \ck in {0,...,#1}{\draw[edge] (0,#1,\ck)--(#1,#1,\ck);}
    \foreach \cj in {0,...,#1}{\draw[edge] (#1,\cj,0)--(#1,\cj,#1);}
    \foreach \ck in {0,...,#1}{\draw[edge] (#1,0,\ck)--(#1,#1,\ck);}
    \foreach \ci in {0,...,#1}{\foreach \cj in {0,...,#1}{\node[v] at (\ci,\cj,0){};}}
    \foreach \ci in {0,...,#1}{\foreach \ck in {0,...,#1}{\node[v] at (\ci,#1,\ck){};}}
    \foreach \cj in {0,...,#1}{\foreach \ck in {0,...,#1}{\node[v] at (#1,\cj,\ck){};}}
  \end{scope}
}

\newcount\pgPlaced
\newcount\pgNf
\newcount\pgTt
\newcommand{\gpancake}[2]{
  \pgNf=1 \pgTt=1
  \loop\ifnum\pgTt<#1 \advance\pgTt by1\relax \multiply\pgNf by\pgTt \repeat
  \edef\pgNfact{\the\pgNf}
  \pgfmathtruncatemacro{\pgTotm}{#1^#1-1}
  \pgPlaced=0
  \foreach \pgI in {0,...,\pgTotm}{
    \gdef\pgPerm{}\xdef\pgUsed{0}\xdef\pgValid{1}
    \foreach \pos in {1,...,#1}{
      \pgfmathtruncatemacro{\pgDig}{Mod(floor(\pgI/(#1^(#1-\pos))),#1)}
      \pgfmathtruncatemacro{\pgBit}{Mod(floor(\pgUsed/(2^\pgDig)),2)}
      \ifnum\pgBit=1 \xdef\pgValid{0}
      \else
        \pgfmathtruncatemacro{\pgNewUsed}{\pgUsed+2^\pgDig}\xdef\pgUsed{\pgNewUsed}
        \pgfmathtruncatemacro{\pgD}{\pgDig+1}\xdef\pgPerm{\pgPerm\pgD}
      \fi}
    \ifnum\pgValid=1
      \pgfmathsetmacro{\pgAng}{90-(\the\pgPlaced)*360/\pgNfact}
      \node[v,label={\pgAng:{\tiny\pgPerm}}] (p\pgPerm) at (\pgAng:#2){};
      \global\advance\pgPlaced by1\relax
    \fi}
  \foreach \pgI in {0,...,\pgTotm}{
    \gdef\pgPerm{}\xdef\pgUsed{0}\xdef\pgValid{1}
    \foreach \pos in {1,...,#1}{
      \pgfmathtruncatemacro{\pgDig}{Mod(floor(\pgI/(#1^(#1-\pos))),#1)}
      \pgfmathtruncatemacro{\pgBit}{Mod(floor(\pgUsed/(2^\pgDig)),2)}
      \ifnum\pgBit=1 \xdef\pgValid{0}
      \else
        \pgfmathtruncatemacro{\pgNewUsed}{\pgUsed+2^\pgDig}\xdef\pgUsed{\pgNewUsed}
        \pgfmathtruncatemacro{\pgD}{\pgDig+1}\xdef\pgPerm{\pgPerm\pgD}
        \expandafter\xdef\csname pgDig\pos\endcsname{\pgD}
      \fi}
    \ifnum\pgValid=1
      \foreach \k in {2,...,#1}{
        \xdef\pgNb{}
        \foreach \j in {1,...,\k}{\xdef\pgNb{\csname pgDig\j\endcsname\pgNb}}
        \ifnum\k<#1
          \foreach \j in {\the\numexpr\k+1\relax,...,#1}{\xdef\pgNb{\pgNb\csname pgDig\j\endcsname}}
        \fi
        \draw[edge] (p\pgPerm)--(p\pgNb);
      }
    \fi}
}




\newcommand{\gqgon}[2]{
  \pgfmathtruncatemacro{\Nm}{2^#1-1}
  \pgfmathtruncatemacro{\nm}{#1-1}
  \foreach \i in {0,...,\Nm}{
    \pgfmathsetmacro{\ang}{90-\i*360/(\Nm+1)}
    \xdef\lbl{}
    \foreach \k in {\nm,...,0}{
      \pgfmathtruncatemacro{\bit}{mod(floor(\i/(2^\k)),2)}
      \xdef\lbl{\lbl\bit}
    }
    \node[v,label={\ang:\scriptsize $\lbl$}] (q\i) at (\ang:#2){};
  }
  \foreach \i in {0,...,\Nm}{
    \foreach \k in {0,...,\nm}{
      \pgfmathtruncatemacro{\bit}{mod(floor(\i/(2^\k)),2)}
      \ifnum\bit=0
        \pgfmathtruncatemacro{\j}{\i+2^\k}
        \draw[edge] (q\i)--(q\j);
      \fi
    }
  }
}




\newcommand{\newinstructor}[2]{
  \definecolor{#2color}{HTML}{#1}
  \expandafter\newcommand\csname #2\endcsname[1]{
    {\color{#2color}\textbf{#2} : ##1}}}

\newcounter{graphnum}
\newif\ifgfbare
\def\gfEmpty{}
\def\gfNUM{NUM}
\newenvironment{graphfig}[2][1]{
  \def\gftmp{#2}
  \ifx\gftmp\gfEmpty
    \gfbaretrue
  \else
    \gfbarefalse
    \ifx\gftmp\gfNUM
      \stepcounter{graphnum}\edef\gflabel{\arabic{graphnum}}
    \else
      \def\gflabel{#2}
    \fi
  \fi
  \ifgfbare\else\begin{center}\fi
  \begin{tikzpicture}[scale=#1]
}{
  \end{tikzpicture}
  \ifgfbare\else\par\vspace{3pt}\textbf{[Graph \gflabel]}\end{center}\fi
}

\newcommand{\glabel}[1]{\begin{center}\textbf{[Graph #1]}\end{center}} (after \documentclass).
%
%  Expects (optionally) these to be defined in main.tex BEFORE \input:
%    \teamname    \reportdate    \logopath
%  Fallbacks are provided below, so nothing breaks if you forget.
%
%  ---------------------------------------------------------------
%  GRAPH TOOLKIT CHEAT SHEET  (all calls go INSIDE graphfig/tikzpicture)
%  ---------------------------------------------------------------
%    \gcycle{n}{radius}{prefix}      n vertices on a circle, joined as a cycle
%    \gcyclen{n}{radius}{prefix}     same, vertices labelled 1..n
%    \gcomplete{n}{radius}{prefix}   complete graph K_n on a circle
%    \gbipartite{m}{n}{gap}          complete bipartite K_{m,n}
%                                      handles: left bl1..blm, right br1..brn
%    \gpath{n}{spacing}{prefix}      path on n vertices, left to right
%    \gpetersen                      Petersen graph  (\gpeterson kept as alias)
%                                      handles: outer po0..po4, inner pin0..pin4
%    \gcube{n}                       n x n x n grid cube, oblique 2D projection
%    \gqgon{n}, \gpancake{n}         suprise hidden functions, only use after you 
%                                      learned what it is. 
%
%  Named handles stay live for the rest of the picture, so after a call:
%    \draw[edge,densely dashed] (po0)--(pin2);   or   \node[hv] at (bl1){};
%
%  Vertex styles:  v (filled dot) . ov (open ring) . hv (highlighted hub)
%  Edge style:     edge
%
%  Figure wrapper: \begin{graphfig}[scale]{Label} ... \end{graphfig}
%    {Text} -> centred, titled  [Graph Text]
%    {NUM}  -> centred, auto-numbered  [Graph 1], [Graph 2], ...
%    {}     -> bare, no title, no centring (use to line pictures up side by side)
%=====================================================================

\usepackage[margin=1in,headheight=30pt,headsep=16pt]{geometry}
\usepackage{amsmath,amssymb,amsfonts,amsthm}
\usepackage{enumitem}
\usepackage{fancyhdr}
\usepackage{graphicx}
\usepackage{tikz}
\usepackage{xcolor}
\usepackage{hyperref}
\usetikzlibrary{arrows.meta,calc}


\providecommand{\teamname}{[GROUP NAME]}
\providecommand{\reportdate}{\today}
\providecommand{\logopath}{assets/logo.png}

\pagestyle{fancy}
\fancyhf{}
\renewcommand{\headrulewidth}{0.4pt}
\renewcommand{\footrulewidth}{0pt}
\fancyhead[L]{\raisebox{2pt}{\small\textsc{\teamname}}}
\fancyhead[R]{\expandafter\IfFileExists\expandafter{\logopath}
  {\includegraphics[height=10pt]{\logopath}}{}}
\fancyfoot[C]{\small\thepage}

\fancypagestyle{firstpage}{
  \fancyhf{}
  \renewcommand{\headrulewidth}{0pt}
  \fancyfoot[C]{\small\thepage}
}

\setlength{\parindent}{0pt}
\setlength{\parskip}{6pt}

\newcommand{\divider}{\begin{center}\rule{0.6\textwidth}{0.4pt}\end{center}}

\newtheoremstyle{solstyle}{6pt}{6pt}{}{}{\bfseries}{.}{ }{}
\theoremstyle{solstyle}
\newtheorem*{solution}{Solution}

\newcommand{\WLoG}{WLoG}
\newcommand{\Cn}[1]{C_{#1}}
\newcommand{\Kn}[1]{K_{#1}}
\newcommand{\Kt}[2]{K_{#1,#2}}

\tikzset{
  v/.style ={circle, fill=black, inner sep=1.6pt},
  ov/.style={circle, draw=black, fill=white, line width=0.6pt, inner sep=1.6pt},
  hv/.style={circle, fill=black!55, draw=black, line width=0.6pt, inner sep=2.6pt},
  edge/.style={line width=0.6pt},
}


\newcommand{\gcycle}[3]{
  \foreach \gi in {1,...,#1}{\node[v] (#3\gi) at ({90-(\gi-1)*360/#1}:#2){};}
  \foreach \gi in {1,...,#1}{\pgfmathtruncatemacro{\gj}{mod(\gi,#1)+1}\draw[edge] (#3\gi)--(#3\gj);}
}

\newcommand{\gcyclen}[3]{
  \foreach \gi in {1,...,#1}{
    \pgfmathsetmacro{\gang}{90-(\gi-1)*360/#1}
    \node[v,label={\gang:\gi}] (#3\gi) at (\gang:#2){};}
  \foreach \gi in {1,...,#1}{\pgfmathtruncatemacro{\gj}{mod(\gi,#1)+1}\draw[edge] (#3\gi)--(#3\gj);}
}

\newcommand{\gcomplete}[3]{
  \foreach \gi in {1,...,#1}{\node[v] (#3\gi) at ({90-(\gi-1)*360/#1}:#2){};}
  \foreach \gi in {1,...,#1}{\foreach \gj in {1,...,#1}{\ifnum\gi<\gj\relax\draw[edge] (#3\gi)--(#3\gj);\fi}}
}

\newcommand{\gbipartite}[3]{
  \foreach \gi in {1,...,#1}{\node[v] (bl\gi) at (0,{(#1+1)/2-\gi}){};}
  \foreach \gj in {1,...,#2}{\node[v] (br\gj) at (#3,{(#2+1)/2-\gj}){};}
  \foreach \gi in {1,...,#1}{\foreach \gj in {1,...,#2}{\draw[edge] (bl\gi)--(br\gj);}}
}

\newcommand{\gpath}[3]{
  \foreach \gi in {1,...,#1}{\node[v] (#3\gi) at ({(\gi-1)*#2},0){};}
  \pgfmathtruncatemacro{\gplast}{#1-1}
  \foreach \gi in {1,...,\gplast}{\pgfmathtruncatemacro{\gpk}{\gi+1}\draw[edge] (#3\gi)--(#3\gpk);}
}


\newcommand{\gpetersen}{
  \foreach \gp in {0,...,4}{
    \node[v] (po\gp) at ({90+\gp*72}:2){};
    \node[v] (pin\gp) at ({90+\gp*72}:1){};}
  \foreach \gp in {0,...,4}{
    \pgfmathtruncatemacro{\gpj}{mod(\gp+1,5)}
    \pgfmathtruncatemacro{\gpk}{mod(\gp+2,5)}
    \draw[edge] (po\gp)--(po\gpj);
    \draw[edge] (po\gp)--(pin\gp);
    \draw[edge] (pin\gp)--(pin\gpk);
  }
}
\newcommand{\gpeterson}{\gpetersen}

\newcommand{\gcube}[1]{
  \begin{scope}[x={(1cm,0cm)}, y={(0cm,1cm)}, z={(0.5cm,0.45cm)}]
    \foreach \ci in {0,...,#1}{\draw[edge] (\ci,0,0)--(\ci,#1,0);}
    \foreach \cj in {0,...,#1}{\draw[edge] (0,\cj,0)--(#1,\cj,0);}
    \foreach \ci in {0,...,#1}{\draw[edge] (\ci,#1,0)--(\ci,#1,#1);}
    \foreach \ck in {0,...,#1}{\draw[edge] (0,#1,\ck)--(#1,#1,\ck);}
    \foreach \cj in {0,...,#1}{\draw[edge] (#1,\cj,0)--(#1,\cj,#1);}
    \foreach \ck in {0,...,#1}{\draw[edge] (#1,0,\ck)--(#1,#1,\ck);}
    \foreach \ci in {0,...,#1}{\foreach \cj in {0,...,#1}{\node[v] at (\ci,\cj,0){};}}
    \foreach \ci in {0,...,#1}{\foreach \ck in {0,...,#1}{\node[v] at (\ci,#1,\ck){};}}
    \foreach \cj in {0,...,#1}{\foreach \ck in {0,...,#1}{\node[v] at (#1,\cj,\ck){};}}
  \end{scope}
}

\newcount\pgPlaced
\newcount\pgNf
\newcount\pgTt
\newcommand{\gpancake}[2]{
  \pgNf=1 \pgTt=1
  \loop\ifnum\pgTt<#1 \advance\pgTt by1\relax \multiply\pgNf by\pgTt \repeat
  \edef\pgNfact{\the\pgNf}
  \pgfmathtruncatemacro{\pgTotm}{#1^#1-1}
  \pgPlaced=0
  \foreach \pgI in {0,...,\pgTotm}{
    \gdef\pgPerm{}\xdef\pgUsed{0}\xdef\pgValid{1}
    \foreach \pos in {1,...,#1}{
      \pgfmathtruncatemacro{\pgDig}{Mod(floor(\pgI/(#1^(#1-\pos))),#1)}
      \pgfmathtruncatemacro{\pgBit}{Mod(floor(\pgUsed/(2^\pgDig)),2)}
      \ifnum\pgBit=1 \xdef\pgValid{0}
      \else
        \pgfmathtruncatemacro{\pgNewUsed}{\pgUsed+2^\pgDig}\xdef\pgUsed{\pgNewUsed}
        \pgfmathtruncatemacro{\pgD}{\pgDig+1}\xdef\pgPerm{\pgPerm\pgD}
      \fi}
    \ifnum\pgValid=1
      \pgfmathsetmacro{\pgAng}{90-(\the\pgPlaced)*360/\pgNfact}
      \node[v,label={\pgAng:{\tiny\pgPerm}}] (p\pgPerm) at (\pgAng:#2){};
      \global\advance\pgPlaced by1\relax
    \fi}
  \foreach \pgI in {0,...,\pgTotm}{
    \gdef\pgPerm{}\xdef\pgUsed{0}\xdef\pgValid{1}
    \foreach \pos in {1,...,#1}{
      \pgfmathtruncatemacro{\pgDig}{Mod(floor(\pgI/(#1^(#1-\pos))),#1)}
      \pgfmathtruncatemacro{\pgBit}{Mod(floor(\pgUsed/(2^\pgDig)),2)}
      \ifnum\pgBit=1 \xdef\pgValid{0}
      \else
        \pgfmathtruncatemacro{\pgNewUsed}{\pgUsed+2^\pgDig}\xdef\pgUsed{\pgNewUsed}
        \pgfmathtruncatemacro{\pgD}{\pgDig+1}\xdef\pgPerm{\pgPerm\pgD}
        \expandafter\xdef\csname pgDig\pos\endcsname{\pgD}
      \fi}
    \ifnum\pgValid=1
      \foreach \k in {2,...,#1}{
        \xdef\pgNb{}
        \foreach \j in {1,...,\k}{\xdef\pgNb{\csname pgDig\j\endcsname\pgNb}}
        \ifnum\k<#1
          \foreach \j in {\the\numexpr\k+1\relax,...,#1}{\xdef\pgNb{\pgNb\csname pgDig\j\endcsname}}
        \fi
        \draw[edge] (p\pgPerm)--(p\pgNb);
      }
    \fi}
}




\newcommand{\gqgon}[2]{
  \pgfmathtruncatemacro{\Nm}{2^#1-1}
  \pgfmathtruncatemacro{\nm}{#1-1}
  \foreach \i in {0,...,\Nm}{
    \pgfmathsetmacro{\ang}{90-\i*360/(\Nm+1)}
    \xdef\lbl{}
    \foreach \k in {\nm,...,0}{
      \pgfmathtruncatemacro{\bit}{mod(floor(\i/(2^\k)),2)}
      \xdef\lbl{\lbl\bit}
    }
    \node[v,label={\ang:\scriptsize $\lbl$}] (q\i) at (\ang:#2){};
  }
  \foreach \i in {0,...,\Nm}{
    \foreach \k in {0,...,\nm}{
      \pgfmathtruncatemacro{\bit}{mod(floor(\i/(2^\k)),2)}
      \ifnum\bit=0
        \pgfmathtruncatemacro{\j}{\i+2^\k}
        \draw[edge] (q\i)--(q\j);
      \fi
    }
  }
}




\newcommand{\newinstructor}[2]{
  \definecolor{#2color}{HTML}{#1}
  \expandafter\newcommand\csname #2\endcsname[1]{
    {\color{#2color}\textbf{#2} : ##1}}}

\newcounter{graphnum}
\newif\ifgfbare
\def\gfEmpty{}
\def\gfNUM{NUM}
\newenvironment{graphfig}[2][1]{
  \def\gftmp{#2}
  \ifx\gftmp\gfEmpty
    \gfbaretrue
  \else
    \gfbarefalse
    \ifx\gftmp\gfNUM
      \stepcounter{graphnum}\edef\gflabel{\arabic{graphnum}}
    \else
      \def\gflabel{#2}
    \fi
  \fi
  \ifgfbare\else\begin{center}\fi
  \begin{tikzpicture}[scale=#1]
}{
  \end{tikzpicture}
  \ifgfbare\else\par\vspace{3pt}\textbf{[Graph \gflabel]}\end{center}\fi
}

\newcommand{\glabel}[1]{\begin{center}\textbf{[Graph #1]}\end{center}} (after \documentclass).
%
%  Expects (optionally) these to be defined in main.tex BEFORE \input:
%    \teamname    \reportdate    \logopath
%  Fallbacks are provided below, so nothing breaks if you forget.
%
%  ---------------------------------------------------------------
%  GRAPH TOOLKIT CHEAT SHEET  (all calls go INSIDE graphfig/tikzpicture)
%  ---------------------------------------------------------------
%    \gcycle{n}{radius}{prefix}      n vertices on a circle, joined as a cycle
%    \gcyclen{n}{radius}{prefix}     same, vertices labelled 1..n
%    \gcomplete{n}{radius}{prefix}   complete graph K_n on a circle
%    \gbipartite{m}{n}{gap}          complete bipartite K_{m,n}
%                                      handles: left bl1..blm, right br1..brn
%    \gpath{n}{spacing}{prefix}      path on n vertices, left to right
%    \gpetersen                      Petersen graph  (\gpeterson kept as alias)
%                                      handles: outer po0..po4, inner pin0..pin4
%    \gcube{n}                       n x n x n grid cube, oblique 2D projection
%    \gqgon{n}, \gpancake{n}         suprise hidden functions, only use after you 
%                                      learned what it is. 
%
%  Named handles stay live for the rest of the picture, so after a call:
%    \draw[edge,densely dashed] (po0)--(pin2);   or   \node[hv] at (bl1){};
%
%  Vertex styles:  v (filled dot) . ov (open ring) . hv (highlighted hub)
%  Edge style:     edge
%
%  Figure wrapper: \begin{graphfig}[scale]{Label} ... \end{graphfig}
%    {Text} -> centred, titled  [Graph Text]
%    {NUM}  -> centred, auto-numbered  [Graph 1], [Graph 2], ...
%    {}     -> bare, no title, no centring (use to line pictures up side by side)
%=====================================================================

\usepackage[margin=1in,headheight=30pt,headsep=16pt]{geometry}
\usepackage{amsmath,amssymb,amsfonts,amsthm}
\usepackage{enumitem}
\usepackage{fancyhdr}
\usepackage{graphicx}
\usepackage{tikz}
\usepackage{xcolor}
\usepackage{hyperref}
\usetikzlibrary{arrows.meta,calc}


\providecommand{\teamname}{[GROUP NAME]}
\providecommand{\reportdate}{\today}
\providecommand{\logopath}{assets/logo.png}

\pagestyle{fancy}
\fancyhf{}
\renewcommand{\headrulewidth}{0.4pt}
\renewcommand{\footrulewidth}{0pt}
\fancyhead[L]{\raisebox{2pt}{\small\textsc{\teamname}}}
\fancyhead[R]{\expandafter\IfFileExists\expandafter{\logopath}
  {\includegraphics[height=10pt]{\logopath}}{}}
\fancyfoot[C]{\small\thepage}

\fancypagestyle{firstpage}{
  \fancyhf{}
  \renewcommand{\headrulewidth}{0pt}
  \fancyfoot[C]{\small\thepage}
}

\setlength{\parindent}{0pt}
\setlength{\parskip}{6pt}

\newcommand{\divider}{\begin{center}\rule{0.6\textwidth}{0.4pt}\end{center}}

\newtheoremstyle{solstyle}{6pt}{6pt}{}{}{\bfseries}{.}{ }{}
\theoremstyle{solstyle}
\newtheorem*{solution}{Solution}

\newcommand{\WLoG}{WLoG}
\newcommand{\Cn}[1]{C_{#1}}
\newcommand{\Kn}[1]{K_{#1}}
\newcommand{\Kt}[2]{K_{#1,#2}}

\tikzset{
  v/.style ={circle, fill=black, inner sep=1.6pt},
  ov/.style={circle, draw=black, fill=white, line width=0.6pt, inner sep=1.6pt},
  hv/.style={circle, fill=black!55, draw=black, line width=0.6pt, inner sep=2.6pt},
  edge/.style={line width=0.6pt},
}


\newcommand{\gcycle}[3]{
  \foreach \gi in {1,...,#1}{\node[v] (#3\gi) at ({90-(\gi-1)*360/#1}:#2){};}
  \foreach \gi in {1,...,#1}{\pgfmathtruncatemacro{\gj}{mod(\gi,#1)+1}\draw[edge] (#3\gi)--(#3\gj);}
}

\newcommand{\gcyclen}[3]{
  \foreach \gi in {1,...,#1}{
    \pgfmathsetmacro{\gang}{90-(\gi-1)*360/#1}
    \node[v,label={\gang:\gi}] (#3\gi) at (\gang:#2){};}
  \foreach \gi in {1,...,#1}{\pgfmathtruncatemacro{\gj}{mod(\gi,#1)+1}\draw[edge] (#3\gi)--(#3\gj);}
}

\newcommand{\gcomplete}[3]{
  \foreach \gi in {1,...,#1}{\node[v] (#3\gi) at ({90-(\gi-1)*360/#1}:#2){};}
  \foreach \gi in {1,...,#1}{\foreach \gj in {1,...,#1}{\ifnum\gi<\gj\relax\draw[edge] (#3\gi)--(#3\gj);\fi}}
}

\newcommand{\gbipartite}[3]{
  \foreach \gi in {1,...,#1}{\node[v] (bl\gi) at (0,{(#1+1)/2-\gi}){};}
  \foreach \gj in {1,...,#2}{\node[v] (br\gj) at (#3,{(#2+1)/2-\gj}){};}
  \foreach \gi in {1,...,#1}{\foreach \gj in {1,...,#2}{\draw[edge] (bl\gi)--(br\gj);}}
}

\newcommand{\gpath}[3]{
  \foreach \gi in {1,...,#1}{\node[v] (#3\gi) at ({(\gi-1)*#2},0){};}
  \pgfmathtruncatemacro{\gplast}{#1-1}
  \foreach \gi in {1,...,\gplast}{\pgfmathtruncatemacro{\gpk}{\gi+1}\draw[edge] (#3\gi)--(#3\gpk);}
}


\newcommand{\gpetersen}{
  \foreach \gp in {0,...,4}{
    \node[v] (po\gp) at ({90+\gp*72}:2){};
    \node[v] (pin\gp) at ({90+\gp*72}:1){};}
  \foreach \gp in {0,...,4}{
    \pgfmathtruncatemacro{\gpj}{mod(\gp+1,5)}
    \pgfmathtruncatemacro{\gpk}{mod(\gp+2,5)}
    \draw[edge] (po\gp)--(po\gpj);
    \draw[edge] (po\gp)--(pin\gp);
    \draw[edge] (pin\gp)--(pin\gpk);
  }
}
\newcommand{\gpeterson}{\gpetersen}

\newcommand{\gcube}[1]{
  \begin{scope}[x={(1cm,0cm)}, y={(0cm,1cm)}, z={(0.5cm,0.45cm)}]
    \foreach \ci in {0,...,#1}{\draw[edge] (\ci,0,0)--(\ci,#1,0);}
    \foreach \cj in {0,...,#1}{\draw[edge] (0,\cj,0)--(#1,\cj,0);}
    \foreach \ci in {0,...,#1}{\draw[edge] (\ci,#1,0)--(\ci,#1,#1);}
    \foreach \ck in {0,...,#1}{\draw[edge] (0,#1,\ck)--(#1,#1,\ck);}
    \foreach \cj in {0,...,#1}{\draw[edge] (#1,\cj,0)--(#1,\cj,#1);}
    \foreach \ck in {0,...,#1}{\draw[edge] (#1,0,\ck)--(#1,#1,\ck);}
    \foreach \ci in {0,...,#1}{\foreach \cj in {0,...,#1}{\node[v] at (\ci,\cj,0){};}}
    \foreach \ci in {0,...,#1}{\foreach \ck in {0,...,#1}{\node[v] at (\ci,#1,\ck){};}}
    \foreach \cj in {0,...,#1}{\foreach \ck in {0,...,#1}{\node[v] at (#1,\cj,\ck){};}}
  \end{scope}
}

\newcount\pgPlaced
\newcount\pgNf
\newcount\pgTt
\newcommand{\gpancake}[2]{
  \pgNf=1 \pgTt=1
  \loop\ifnum\pgTt<#1 \advance\pgTt by1\relax \multiply\pgNf by\pgTt \repeat
  \edef\pgNfact{\the\pgNf}
  \pgfmathtruncatemacro{\pgTotm}{#1^#1-1}
  \pgPlaced=0
  \foreach \pgI in {0,...,\pgTotm}{
    \gdef\pgPerm{}\xdef\pgUsed{0}\xdef\pgValid{1}
    \foreach \pos in {1,...,#1}{
      \pgfmathtruncatemacro{\pgDig}{Mod(floor(\pgI/(#1^(#1-\pos))),#1)}
      \pgfmathtruncatemacro{\pgBit}{Mod(floor(\pgUsed/(2^\pgDig)),2)}
      \ifnum\pgBit=1 \xdef\pgValid{0}
      \else
        \pgfmathtruncatemacro{\pgNewUsed}{\pgUsed+2^\pgDig}\xdef\pgUsed{\pgNewUsed}
        \pgfmathtruncatemacro{\pgD}{\pgDig+1}\xdef\pgPerm{\pgPerm\pgD}
      \fi}
    \ifnum\pgValid=1
      \pgfmathsetmacro{\pgAng}{90-(\the\pgPlaced)*360/\pgNfact}
      \node[v,label={\pgAng:{\tiny\pgPerm}}] (p\pgPerm) at (\pgAng:#2){};
      \global\advance\pgPlaced by1\relax
    \fi}
  \foreach \pgI in {0,...,\pgTotm}{
    \gdef\pgPerm{}\xdef\pgUsed{0}\xdef\pgValid{1}
    \foreach \pos in {1,...,#1}{
      \pgfmathtruncatemacro{\pgDig}{Mod(floor(\pgI/(#1^(#1-\pos))),#1)}
      \pgfmathtruncatemacro{\pgBit}{Mod(floor(\pgUsed/(2^\pgDig)),2)}
      \ifnum\pgBit=1 \xdef\pgValid{0}
      \else
        \pgfmathtruncatemacro{\pgNewUsed}{\pgUsed+2^\pgDig}\xdef\pgUsed{\pgNewUsed}
        \pgfmathtruncatemacro{\pgD}{\pgDig+1}\xdef\pgPerm{\pgPerm\pgD}
        \expandafter\xdef\csname pgDig\pos\endcsname{\pgD}
      \fi}
    \ifnum\pgValid=1
      \foreach \k in {2,...,#1}{
        \xdef\pgNb{}
        \foreach \j in {1,...,\k}{\xdef\pgNb{\csname pgDig\j\endcsname\pgNb}}
        \ifnum\k<#1
          \foreach \j in {\the\numexpr\k+1\relax,...,#1}{\xdef\pgNb{\pgNb\csname pgDig\j\endcsname}}
        \fi
        \draw[edge] (p\pgPerm)--(p\pgNb);
      }
    \fi}
}




\newcommand{\gqgon}[2]{
  \pgfmathtruncatemacro{\Nm}{2^#1-1}
  \pgfmathtruncatemacro{\nm}{#1-1}
  \foreach \i in {0,...,\Nm}{
    \pgfmathsetmacro{\ang}{90-\i*360/(\Nm+1)}
    \xdef\lbl{}
    \foreach \k in {\nm,...,0}{
      \pgfmathtruncatemacro{\bit}{mod(floor(\i/(2^\k)),2)}
      \xdef\lbl{\lbl\bit}
    }
    \node[v,label={\ang:\scriptsize $\lbl$}] (q\i) at (\ang:#2){};
  }
  \foreach \i in {0,...,\Nm}{
    \foreach \k in {0,...,\nm}{
      \pgfmathtruncatemacro{\bit}{mod(floor(\i/(2^\k)),2)}
      \ifnum\bit=0
        \pgfmathtruncatemacro{\j}{\i+2^\k}
        \draw[edge] (q\i)--(q\j);
      \fi
    }
  }
}




\newcommand{\newinstructor}[2]{
  \definecolor{#2color}{HTML}{#1}
  \expandafter\newcommand\csname #2\endcsname[1]{
    {\color{#2color}\textbf{#2} : ##1}}}

\newcounter{graphnum}
\newif\ifgfbare
\def\gfEmpty{}
\def\gfNUM{NUM}
\newenvironment{graphfig}[2][1]{
  \def\gftmp{#2}
  \ifx\gftmp\gfEmpty
    \gfbaretrue
  \else
    \gfbarefalse
    \ifx\gftmp\gfNUM
      \stepcounter{graphnum}\edef\gflabel{\arabic{graphnum}}
    \else
      \def\gflabel{#2}
    \fi
  \fi
  \ifgfbare\else\begin{center}\fi
  \begin{tikzpicture}[scale=#1]
}{
  \end{tikzpicture}
  \ifgfbare\else\par\vspace{3pt}\textbf{[Graph \gflabel]}\end{center}\fi
}

\newcommand{\glabel}[1]{\begin{center}\textbf{[Graph #1]}\end{center}}